% Supplementary file: full criterion audits.
% Per claim: one page of Source + Description + verdict table + readings table,
% then a new page carrying the full 35-criterion scorecard.
\pdfoutput=1
\documentclass{article}
\usepackage[preprint]{tmlr}
\newcommand{\openreview}{\url{https://openreview.net/forum?id=XXXX}}
\def\month{06}
\def\year{2026}
\usepackage{amsmath,amssymb}
\usepackage{booktabs}
\usepackage{array}
\usepackage{float}
\usepackage{longtable}
\usepackage{xcolor}
\usepackage{caption}
\usepackage{hyperref}   % defines \url/urlstyle, without which \doi{} breaks on underscores
\captionsetup[table]{skip=8pt}
\setlength{\intextsep}{16pt plus 2pt minus 2pt}

\title{Supplementary Material:\\Criterion Audits of Mechanistic Claims}
\author{}

\begin{document}
\maketitle

\section*{How to read these audits}

Each claim occupies two pages: a summary page carrying the verdict and how it was reached,
and a facing page carrying the full 35-criterion scorecard. Three conventions govern the
scoring.

\textbf{Scope is part of the claim.} A mechanistic claim is not true or false on its own;
it is true or false at a scope. We report the broadest scope at which the claim reaches
each verdict. A claim that holds narrowly and fails broadly is recorded as both, because
reporting only the failure misrepresents the origin paper and reporting only the success
misrepresents how the claim is used. The scorecard scores the primary reading; the readings table records the others.

\textbf{The verdict comes from gates, not from counts.} The status tally at the foot of
each scorecard is descriptive. Promotion between tiers is decided by named criteria, and
the verdict table reports which one blocks it.

\textbf{Untested splits two ways.} \emph{Unattempted} means the experiment was not run.
\emph{Untestable in the setting as run} means it could not have been. Only the first is a
gap the authors could have closed.


% ---- Induction Heads: In-Context Token Copying (Triangulated)
\input{generated/induction_section_proto}


% ---- Copy Suppression in Negative Attention Heads (Mechanistically Supported)


% ---- Greater-Than Circuit in Year-Span Prediction (Mechanistically Supported)


% ---- Fourier Multiplication in Modular Addition (Mechanistically Supported)

% ---- Refusal Mediated by a Single Direction (Mechanistically Supported)


% ---- Successor Heads: Ordinal Incrementation (Mechanistically Supported)


% ---- Superposition of Features in Toy Models (Mechanistically Supported)


% ---- Global Workspace / J-space (Mechanistically Supported)


% ---- Docstring Circuit (Causally Suggestive)


% ---- Gender Bias Circuits (Causally Suggestive)


% ---- IOI Circuit (Causally Suggestive)


% ---- Othello World Model (Causally Suggestive)


% ---- SAE Features as Units of Analysis (Causally Suggestive)


% ---- Probing Classifiers as Mechanistic Evidence (Proposed)


% ---- Induction Heads: General In-Context Learning (Disconfirmed)


% ---- Knowledge Neurons (Disconfirmed)


\bibliographystyle{tmlr}
\bibliography{references}

\end{document}
